\chapter{Thiết lập bổ trợ và định nghĩa chỉ số}
\label{app:supporting-methods}

\section{Khảo sát các bộ dữ liệu và cơ sở lựa chọn}
\label{sec:dataset-selection}
Việc lựa chọn dữ liệu xuất phát từ mục tiêu phân loại glaucoma
bằng khối ảnh OCT B-scan và các ảnh chiếu của chính khối
ảnh đó. Các tiêu chí ưu tiên là: có chuỗi B-scan tổ chức thành khối
3D; có nhãn phù hợp với bài toán phân loại; vùng quét phù hợp với
phạm vi nghiên cứu; và có quy mô cùng phân hoạch cho phép huấn luyện,
lựa chọn và đánh giá mô hình. Số ảnh có thể dùng để khử nhiễu không
phải tiêu chí quyết định chọn bộ dữ liệu.

\begingroup
\setlength{\tabcolsep}{5pt}
\begin{longtable}{@{}>{\RaggedRight\arraybackslash}p{2.2cm}>{\RaggedRight\arraybackslash}p{3.3cm}>{\RaggedRight\arraybackslash}p{2.9cm}>{\RaggedRight\arraybackslash}p{6.1cm}@{}}
  \caption{Đối chiếu các bộ dữ liệu với bài toán phân loại từ khối B-scan và ảnh chiếu}
  \label{tab:dataset-comparison}\\
  \toprule
  \textbf{Bộ dữ liệu} & \textbf{Dữ liệu OCT và vùng quét} &
  \textbf{Quy mô công bố} & \textbf{Nhãn và mức phù hợp với luận văn}\\
  \midrule
  \endfirsthead
  \multicolumn{4}{c}{\textit{Bảng~\thetable{} (tiếp theo)}}\\
  \toprule
  \textbf{Bộ dữ liệu} & \textbf{Dữ liệu OCT và vùng quét} &
  \textbf{Quy mô công bố} & \textbf{Nhãn và mức phù hợp với luận văn}\\
  \midrule
  \endhead
  \bottomrule
  \endfoot
  Harvard-GF \cite{luo2024harvardgf,harvardgfRepo} &
  Khối B-scan $200^3$; vùng đĩa thị &
  3.300 khối ảnh từ 3.300 người &
  Có nhãn âm tính/dương tính. Chọn làm nguồn chính cho nhánh 3D
  và các ảnh chiếu 2D.\\
  \midrule
  GAMMA \cite{wu2022gamma} &
  Khối OCT vùng hoàng điểm; 256 B-scan/khối; kèm ảnh đáy mắt &
  300 cặp OCT--đáy mắt từ 276 người &
  Có nhãn phân độ bệnh; thử thách chính gộp thành ba mức.
  Là ứng viên khảo sát khả năng khái quát, nhưng khác vùng quét
  và cần thống nhất nhãn trước khi đánh giá ngoài.\\
  \midrule
  OCTDL \cite{kulyabin2024octdl} &
  B-scan 2D vùng hoàng điểm, không phát hành dưới dạng khối 3D &
  2.064 ảnh từ 821 người &
  Bảy nhóm bệnh/trạng thái võng mạc, không có lớp bệnh đích.
  Không đáp ứng đồng thời đầu vào 3D và nhãn phân loại cần thiết.\\
  \midrule
  OCT2017 / Kermany \cite{kermany2018data} &
  B-scan 2D; các ảnh được chia theo nhóm nhãn &
  84.484 ảnh theo thống kê tác giả cung cấp\textsuperscript{*} &
  Bốn nhóm CNV, DME, drusen và bình thường; không có lớp bệnh đích.
  Số ảnh 2D lớn không thay thế số khối 3D có nhãn phù hợp.\\
  \midrule
  OCTID \cite{gholami2018octid} &
  Ảnh OCT 2D từ quét raster võng mạc &
  Hơn 500 ảnh &
  Năm nhóm NO, MH, AMD, CSR và DR; không có lớp bệnh đích.
  Không chọn làm dữ liệu huấn luyện chính.\\
  \midrule
  Duke / Srinivasan \cite{srinivasan2014data} &
  Khối OCT vùng hoàng điểm &
  45 người: 15 bình thường, 15 AMD, 15 DME &
  Có dữ liệu 3D nhưng không có nhãn bệnh đích và quy mô nhỏ.
  Không dùng để đánh giá trực tiếp bài toán phân loại của luận văn.\\
  \midrule
  Harvard-GDP \cite{harvardgdpRepo} &
  Bản đồ RNFLT $225\times225$ và thị trường trong bản phát hành &
  1.000 người &
  Có nhãn bệnh và nhãn tiến triển cho một phần mẫu, nhưng không có
  trường khối B-scan trong mô tả bản phát hành; không phù hợp với
  nguồn đầu vào đã giới hạn của luận văn.\\
\end{longtable}
\endgroup
\noindent\textit{Ghi chú:} Quy mô được ghi theo đơn vị của từng nguồn,
không so sánh trực tiếp số ảnh 2D với số khối 3D hoặc số người.
\textsuperscript{*}Số ảnh OCT2017 cần đối chiếu lại với danh sách tệp
của phiên bản tải về; không mặc định mọi bản sao có cùng số ảnh.

Từ đối chiếu trên, luận văn chọn Harvard-GF vì bộ dữ liệu cung cấp
đồng thời khối B-scan 3D vùng đĩa thị và nhãn phân loại nhị phân.
Nhờ đó, nhánh 3D và các nhánh ảnh chiếu có thể được xây dựng từ
cùng một mẫu, tạo cơ sở so sánh mô hình đơn nhánh với mô hình hợp
nhất mà không cần bổ sung loại dữ liệu lâm sàng khác. Quy mô 3.300
mẫu và phân hoạch có sẵn thuận tiện cho thiết kế thực nghiệm,
nhưng không bảo đảm một kiến trúc 3D cụ thể sẽ hội tụ hoặc khái quát tốt.

Đối với GAMMA, khác biệt về vùng quét và cách định nghĩa nhãn là
hạn chế chính khi áp dụng trực tiếp thiết kế hiện tại. Việc đánh
giá ngoài trên bộ dữ liệu này đòi hỏi xác định trước
ánh xạ nhãn, kiểm tra tính tương thích của phép chiếu và phân tích
sự khác biệt phân phối \cite{wu2022gamma}. Các bộ OCT 2D có thể
được xem xét cho nghiên cứu chuyển giao riêng, nhưng không nằm trong
quy trình huấn luyện của luận văn. Khả năng hỗ trợ khử nhiễu không
được dùng để thay thế yêu cầu về nhãn và cấu trúc dữ liệu phân loại.

\section{Hệ thống chỉ số và quy ước đánh giá}
\label{sec:evaluation-metrics}
Các chỉ số được tổ chức theo đối tượng đo trong Bảng~\ref{tab:variables}.
Val và Test lần lượt chỉ tập xác thực và tập kiểm tra, không phải tên chỉ số.
Chẳng hạn, Val AUROC khác Val accuracy; riêng hai cột Val và Test của khảo
sát nền ở Mục~\ref{sec:app-3d-baseline} được dùng với nghĩa độ chính xác.
Các công thức dưới đây xác định cách diễn giải và quy trình báo cáo.
Đối với số liệu đã có, những quy ước chưa được lưu trong nhật ký vẫn cần
đối chiếu mã đánh giá; không mặc định mọi lần chạy đã dùng cùng cách tính.

\begin{table}[htbp]
  \centering
  \caption{Các nhóm chỉ số trong luận văn và phạm vi sử dụng}
  \label{tab:variables}
  \begin{tabular}{@{}>{\RaggedRight\arraybackslash}p{0.22\textwidth}>{\RaggedRight\arraybackslash}p{0.39\textwidth}>{\RaggedRight\arraybackslash}p{0.30\textwidth}@{}}
    \toprule
    Nhóm & Đại lượng & Phạm vi\\ \midrule
    Nhãn phân loại & Accuracy, precision, recall, specificity, F1, balanced accuracy, MCC & Đánh giá tại ngưỡng cố định\\
    Điểm số và xác suất & AUROC, PR-AUC, loss, ECE & Phân biệt và hiệu chỉnh xác suất\\
    Quá trình học và tài nguyên & bestEp, best/final, mức giảm; tham số, VRAM, thời gian & Nhật ký huấn luyện\\
    Chất lượng ảnh & SNR, ENL, CNR, $\beta$ & Khảo sát khử nhiễu bổ trợ\\
    Độ tin cậy & Trung bình, độ lệch chuẩn, khoảng tin cậy & Chưa có thống kê nhiều seed; bootstrap dự đoán cố định nếu đủ dữ liệu\\
    \bottomrule
  \end{tabular}
\end{table}

\subsection{Đơn vị đánh giá và ma trận nhầm lẫn}
Xét tập đánh giá gồm $n$ quan sát, với nhãn thật $y_i\in\{0,1\}$,
trong đó 1 là lớp dương tính, và xác suất dự đoán $p_i=P_\theta(Y=1\mid X_i)$.
Nhãn dự đoán là $\widehat y_i=\mathbf 1\{p_i\geq\tau\}$.
Ngưỡng $\tau$ được xác định trên tập phát triển và giữ nguyên khi kiểm tra.
Một quan sát ứng với một khối OCT theo phân hoạch đã công bố, không phải
mỗi B-scan hoặc mỗi ảnh chiếu của cùng khối. Nếu giữ nhiều mắt hoặc lần
chụp của một người, phân hoạch và ước lượng sai số phải xét phụ thuộc theo
bệnh nhân. Các số đếm được cộng trên toàn tập trước khi tính chỉ số.

\begin{equation}
\begin{aligned}
 TP&=\sum_{i=1}^{n}\mathbf 1\{y_i=1,\widehat y_i=1\},&
 FN&=\sum_{i=1}^{n}\mathbf 1\{y_i=1,\widehat y_i=0\},\\
 FP&=\sum_{i=1}^{n}\mathbf 1\{y_i=0,\widehat y_i=1\},&
 TN&=\sum_{i=1}^{n}\mathbf 1\{y_i=0,\widehat y_i=0\}.
\end{aligned}
\label{eq:metric-confusion}
\end{equation}
TP và TN là số dự đoán đúng tương ứng ở lớp dương và âm; FN là số ca
dương bị bỏ sót, còn FP là số ca âm bị báo dương. Tổng bốn số bằng $n$.
Khi trình bày ma trận nhầm lẫn, hàng biểu thị nhãn thật, cột biểu thị nhãn
dự đoán và thứ tự lớp là $(0,1)$, tức ma trận
$\left(\begin{smallmatrix}TN&FP\\FN&TP\end{smallmatrix}\right)$.

\subsection{Các chỉ số tính từ nhãn dự đoán}
Các công thức nhị phân sau dùng lớp dương tính làm lớp quan tâm
\cite{sklearnMetricsGuide}. Trong đó, recall cũng được gọi là sensitivity
hoặc TPR, còn specificity tương ứng với TNR.
\begingroup
\setlength{\jot}{9pt}
% Hai vế có hộp rộng 0 để dấu bằng nằm đúng giữa vùng văn bản.
\newcommand{\CanGiuaDauBang}[2]{%
  \makebox[0pt][r]{$\displaystyle #1$}
  = \makebox[0pt][l]{$\displaystyle #2$}%
}
\begin{gather}
  \CanGiuaDauBang{\mathrm{Accuracy}}{\frac{TP+TN}{n},}
  \label{eq:metric-acc-precision}\\
  \CanGiuaDauBang{\mathrm{Precision}}{\frac{TP}{TP+FP},}
  \label{eq:metric-precision}\\
  \CanGiuaDauBang{\mathrm{Recall}}{\frac{TP}{TP+FN},}
  \label{eq:metric-recall-spec}\\
  \CanGiuaDauBang{\mathrm{Specificity}}{\frac{TN}{TN+FP},}\\
  \CanGiuaDauBang{\mathrm{FPR}}{\frac{FP}{FP+TN}=1-\mathrm{Specificity},}\\
  \CanGiuaDauBang{\mathrm{Balanced\ accuracy}}{\frac{\mathrm{TPR}+\mathrm{TNR}}{2},}
  \label{eq:metric-balanced}\\
  \CanGiuaDauBang{F1}{\frac{2TP}{2TP+FP+FN}.}
  \label{eq:metric-f1}
\end{gather}
\endgroup
Accuracy là tỷ lệ dự đoán đúng trên toàn tập. Precision là tỷ lệ dự đoán
dương đúng trong tổng số dự đoán dương; recall là tỷ lệ ca dương thật được phát hiện.
Specificity đo khả năng nhận đúng lớp âm, còn FPR đo tỷ lệ báo dương sai
trong lớp âm. Balanced accuracy coi trọng hai lớp như nhau. F1 kết hợp
precision và recall, nhưng không sử dụng TN. Các chỉ số nằm trong $[0,1]$
khi mẫu số khác 0; giá trị lớn hơn được ưu tiên, ngoại trừ FPR.

Hệ số tương quan Matthews sử dụng cả bốn số đếm \cite{sklearnMCC}:
\begin{equation}
 \mathrm{MCC}=\frac{TP\,TN-FP\,FN}
 {\sqrt{(TP+FP)(TP+FN)(TN+FP)(TN+FN)}}.
 \label{eq:metric-mcc}
\end{equation}
MCC nằm trong $[-1,1]$: 1 ứng với dự đoán hoàn toàn đúng, $-1$ ứng với
đảo ngược hoàn toàn và 0 biểu thị không có tương quan tuyến tính giữa hai
nhãn nhị phân khi công thức xác định. Với dự đoán một lớp, mẫu số có thể
bằng 0; số 0 do thư viện trả về là quy ước xử lý, không phải phép chia hợp lệ.

Việc diễn giải trường hợp mẫu số bằng 0 phụ thuộc quy ước của mã đánh
giá. Các công thức trên không thay thế cho trung bình macro hoặc weighted
nếu nhật ký dùng các chế độ đó. Trung bình F1 từng batch không tương đương F1
toàn tập. Accuracy hoặc F1 cao riêng lẻ không loại trừ dự đoán một lớp;
ví dụ cụ thể của các lần chạy hợp nhất đặc trưng được phân tích tại
Phương trình~\eqref{eq:fusion-one-class}.

\subsection{Diện tích dưới đường cong ROC và precision--recall}
AUROC sử dụng điểm số chưa áp ngưỡng, không dùng nhãn đã làm tròn
\cite{sklearnROC}. Đường ROC biểu diễn TPR theo FPR khi thay đổi ngưỡng.
Với $n_+$ mẫu dương và $n_-$ mẫu âm, cách tính thực nghiệm tương đương là
\begin{equation}
 \widehat{\mathrm{AUROC}}=\frac{1}{n_+n_-}
 \sum_{i:y_i=1}\sum_{j:y_j=0}
 \left[\mathbf 1\{p_i>p_j\}+\frac12\mathbf 1\{p_i=p_j\}\right].
 \label{eq:metric-auroc}
\end{equation}
Chỉ số phản ánh khả năng xếp điểm mẫu dương cao hơn mẫu âm, với trọng số một nửa khi hòa điểm.
AUROC bằng 1 là phân biệt hoàn hảo; điểm số hằng cho AUROC bằng 0,5 nếu
cả hai lớp đều có mặt. Giá trị lớn hơn được ưu tiên, nhưng AUROC trên 0,5
trong một lần chạy chưa xác lập ý nghĩa thống kê hoặc độ tin cậy xác suất.

Đường precision--recall biểu diễn precision theo recall. Với các điểm
$(R_k,P_k)$ được sắp theo recall tăng dần, hai cách tổng hợp cần phân biệt là
\begin{align}
 \mathrm{PR\!\text{-}\!AUC}_{\mathrm{trap}}
   &=\sum_{k=1}^{K}(R_k-R_{k-1})\frac{P_k+P_{k-1}}{2},\\
 \mathrm{AP}&=\sum_{k=1}^{K}(R_k-R_{k-1})P_k.
 \label{eq:metric-pr-ap}
\end{align}
AP là average precision, gán trọng số theo phần recall tăng thêm và không
nội suy tuyến tính như quy tắc hình thang \cite{sklearnAP}.
Cả hai đều ưu tiên giá trị lớn hơn nhưng có thể cho kết quả khác nhau.
Cách dựng điểm đầu/cuối và xử lý điểm số hòa là thông tin cần thiết để tái lập.
Luận văn giữ tên PR-AUC của nhật ký
cho đến khi xác nhận hàm tính, thay vì diễn giải các giá trị đó là AP.
Tỷ lệ lớp dương $\widehat\pi=n_+/n$ cần được báo cáo kèm theo để diễn giải
precision và so sánh các tập có phân bố nhãn khác nhau.

Độ nhạy tại mức đặc hiệu mục tiêu $s_0$ được đánh giá bằng cách chọn
$\tau^*$ trên tập xác thực theo quy tắc định trước, chẳng hạn
\begin{equation}
 \tau^*\in\underset{\tau:\,\mathrm{Specificity}_{\mathrm{val}}(\tau)\geq s_0}
 {\operatorname{argmax}}\ \mathrm{Recall}_{\mathrm{val}}(\tau).
 \label{eq:metric-target-specificity}
\end{equation}
Sau khi xử lý hòa điểm theo quy tắc định trước, giữ nguyên $\tau^*$ để tính
recall và specificity trên tập kiểm tra. Độ đặc hiệu thực tế trên tập kiểm tra không mặc nhiên
bằng hoặc vượt $s_0$; cả hai kết quả phải được báo cáo. Chỉ số này còn
thuộc kế hoạch đánh giá, chưa có giá trị trong các bảng hiện tại.

\subsection{Hàm mất mát và chất lượng xác suất dự đoán}
Với xác suất lớp dương và không dùng trọng số, cross-entropy nhị phân
(BCE), cũng gọi là log-loss, được tính bằng \cite{sklearnLogLoss}
\begin{equation}
 L_{\mathrm{BCE}}=-\frac1n\sum_{i=1}^{n}
 \left[y_i\ln \widetilde p_i+(1-y_i)\ln(1-\widetilde p_i)\right],
 \quad \widetilde p_i=\min(1-\varepsilon,\max(\varepsilon,p_i)).
 \label{eq:metric-bce}
\end{equation}
Ở đây $0<\varepsilon<1/2$ tránh logarithm của 0 và phải được ghi cùng
kiểu dữ liệu số. Với hai logits $(a_{i0},a_{i1})$, xác suất tương ứng là
$p_i=\exp(a_{i1})/[\exp(a_{i0})+\exp(a_{i1})]$; triển khai trên logits
cần dạng ổn định số. Loss nhỏ hơn là tốt hơn; dự đoán sai nhưng gán xác
suất rất cao bị phạt mạnh. Loss không đồng nhất với tỷ lệ phân loại sai.
Các lần chạy có trọng số lớp, label smoothing hoặc cách lấy trung bình khác
cần được xác định công thức thực dùng trước khi so sánh. Không mặc định cột loss
đã lưu chính là BCE không trọng số. Logarithm tự nhiên cho đơn vị nat;
đổi sang bit bằng $L_{\mathrm{bit}}=L_{\mathrm{BCE}}/\ln2$.

Sai số hiệu chỉnh kỳ vọng (ECE) đo chênh lệch giữa mức tự tin và tỷ lệ
đúng theo các khoảng xác suất \cite{guo2017calibration}. Với định nghĩa
top-label, đặt $\widehat y_i^{\mathrm{top}}=\mathbf 1\{p_i\geq1/2\}$ và
$c_i=\max(p_i,1-p_i)$. Chia $c_i$ thành $M$ khoảng, tập chỉ số trong
khoảng thứ $m$ là $B_m$. Trên các khoảng không rỗng,
\begin{align}
 a_m&=\frac{1}{|B_m|}\sum_{i\in B_m}\mathbf 1\{y_i=\widehat y_i^{\mathrm{top}}\},&
 c_m&=\frac{1}{|B_m|}\sum_{i\in B_m}c_i,\\
 \mathrm{ECE}_{\mathrm{top}}&=\sum_{m:|B_m|>0}\frac{|B_m|}{n}|a_m-c_m|.
 \label{eq:metric-ece}
\end{align}
ECE thuộc $[0,1]$, càng nhỏ càng tốt theo cách chia khoảng đã chọn.
Một biến thể khác chia khoảng theo $p_i$ và thay $a_m,c_m$ bằng trung
bình $y_i,p_i$ trong khoảng, nhằm đánh giá xác suất lớp dương. Hai định
nghĩa không được dùng lẫn nhau. Việc diễn giải ECE đã ghi nhận đòi hỏi
xác nhận biến thể, số khoảng, biên khoảng và cách chia đều độ rộng hay số mẫu.
ECE thấp không bảo đảm phân biệt tốt; trị tuyệt đối cũng không cho biết
riêng chiều quá tự tin hay thiếu tự tin.

\subsection{Các mốc lựa chọn checkpoint}
Gọi $m_e$ là chỉ số xác thực cần tối đa hóa tại epoch $e$, $E$ là epoch
cuối được ghi. Khi quy định chọn epoch sớm nhất nếu hòa điểm,
\begin{equation}
 e^*=\min\operatorname*{argmax}_{1\leq e\leq E}m_e,\qquad
 m_{\mathrm{best}}=m_{e^*},\quad m_{\mathrm{final}}=m_E,\quad
 \Delta m=m_{\mathrm{best}}-m_{\mathrm{final}}.
 \label{eq:metric-checkpoint}
\end{equation}
bestEp là $e^*$, không phải một chỉ số chất lượng cần tối đa hóa.
Min loss là $\min_e L_e$ và có thể thuộc epoch khác $e^*$.
Mức giảm $\Delta m$ chỉ mô tả chênh lệch hai mốc, không đo đầy đủ dao
động đường học hoặc độ ổn định giữa các hạt giống ngẫu nhiên. Quy tắc hòa điểm trên là quy
ước cho thiết kế đánh giá; các lần chạy trước vẫn cần đối chiếu mã checkpoint.
Khảo sát nền dùng accuracy xác thực, bảng hợp nhất đặc trưng dùng AUROC xác thực
tốt nhất, còn bảng 2D hiện ghi chỉ số xác thực ở epoch cuối.

\subsection{Số tham số, bộ nhớ và thời gian tính toán}
Với các tensor tham số phân biệt $\theta_j$ có kích thước $d_{j1},\ldots,d_{jq_j}$,
số tham số tổng và số tham số được cập nhật là
\begin{equation}
 N_{\mathrm{param}}=\sum_j\prod_{k=1}^{q_j}d_{jk},\qquad
 N_{\mathrm{train}}=\sum_{j:\,\theta_j\ \text{được tối ưu}}\prod_{k=1}^{q_j}d_{jk}.
 \label{eq:metric-parameters}
\end{equation}
Tham số dùng chung chỉ đếm một lần; cột triệu tham số là
$N_{\mathrm{param}}/10^6$. Ít tham số hơn biểu thị mô hình gọn hơn,
không tự bảo đảm nhanh hơn hoặc chính xác hơn.

Bộ nhớ cực đại trên khoảng đo $\mathcal T$ là
$V_{\max}=\max_{t\in\mathcal T}V(t)$. Cần phân biệt bộ nhớ tensor đã
cấp phát với bộ nhớ bộ cấp phát GPU giữ lại; phép đo cực đại phải đặt
lại trước khoảng đo \cite{pytorchMemoryMetric}. Báo cáo rõ GB
($10^9$ byte) hay GiB ($2^{30}$ byte), batch size, kích thước đầu vào,
độ chính xác số và có bao gồm bước lan truyền ngược hay không.
Các cột VRAM của khảo sát trước được giữ nguyên đơn vị do thiếu siêu dữ liệu để xác nhận phép chuyển đổi.

Gọi $t_e$ là thời gian epoch thứ $e$ trong tập epoch được đo
$\mathcal E$, $T_{\mathrm{infer}}$ là thời gian xử lý $n_{\mathrm{infer}}$ mẫu:
\begin{equation}
 \overline t_{\mathrm{ep}}=\frac{\sum_{e\in\mathcal E}t_e}{|\mathcal E|},\qquad
 \ell=\frac{T_{\mathrm{infer}}}{n_{\mathrm{infer}}},\qquad
 r=\frac{n_{\mathrm{infer}}}{T_{\mathrm{infer}}}.
 \label{eq:metric-timing}
\end{equation}
Cột s/ep là giây mỗi epoch; $\ell$ là thời gian trung bình trên mẫu
trong chế độ batch đã đo, còn $r$ là thông lượng (mẫu/giây).
Thời gian toàn bộ lần chạy là $t_{\mathrm{end}}-t_{\mathrm{start}}$, chia 60 hoặc
3600 để đổi giây sang phút hoặc giờ. Thời gian khử nhiễu giây/khối ảnh
cũng phải ghi phạm vi I/O. Khi đo GPU cần đồng bộ trước và sau khoảng
đo, quy định warm-up và cùng điều kiện phần cứng. Thời gian toàn bộ lần chạy
khác thời gian suy luận; các lần chạy khác số epoch không tạo thành phép so sánh chi phí ngang bằng.

\subsection{Độ biến thiên và khoảng tin cậy}
Với $K\geq2$ lần chạy cùng cấu hình, kết quả $m_1,\ldots,m_K$ được tổng hợp bởi
\begin{equation}
 \overline m=\frac1K\sum_{k=1}^{K}m_k,\qquad
 s_m=\sqrt{\frac{1}{K-1}\sum_{k=1}^{K}(m_k-\overline m)^2}.
 \label{eq:metric-seed-statistics}
\end{equation}
$\overline m\pm s_m$ mô tả biến thiên giữa các hạt giống ngẫu nhiên, không phải khoảng tin cậy
95\%. Sáu phương pháp khác nhau không tương đương sáu lần chạy của một cấu
hình; một lần chạy không đủ để tính $s_m$.
Kế hoạch khoảng tin cậy dùng bootstrap theo bệnh nhân: lấy mẫu lại bệnh
nhân có hoàn lại, giữ các quan sát của người được chọn và tính lại chỉ số.
Với các giá trị $m^{*(1)},\ldots,m^{*(B)}$, khoảng percentile 95\% là \cite{efron1986bootstrap}
\begin{equation}
 \mathrm{CI}_{95\%}=\left[Q_{0{,}025}(m^*),Q_{0{,}975}(m^*)\right],
 \label{eq:metric-bootstrap}
\end{equation}
trong đó $Q_q$ là phân vị mức $q$ của phân bố bootstrap thực nghiệm.
Quy trình báo cáo cần nêu $B$, cách xử lý mẫu lại thiếu một lớp và đơn vị đánh giá.
So sánh ghép cặp dùng cùng mẫu lại cho hai mô hình rồi tính chênh lệch.
Khoảng này đánh giá bất định lấy mẫu với mô hình cố định, không thay thế
chạy lặp huấn luyện. Dữ liệu hiện có của luận văn không đủ để tính các khoảng này.

\subsection{Chỉ số chất lượng ảnh}
SNR, ENL, CNR và tương quan gradient $\beta$ được định nghĩa tại
Mục~\ref{sec:denoise-metric-definitions}, gắn với vùng đo của Mục~\ref{sec:denoise-comparison-method}.
Chúng mô tả biến đổi ảnh, không thay thế đánh giá phân loại.

\section{NLM trong đối chứng tiền xử lý}
Phần này cung cấp cơ sở cho bước tiền xử lý của quy trình đa nhánh.
Việc lựa chọn bộ lọc cần được đối chứng trong bài toán phân loại.

\subsection{Nhiễu speckle trong ảnh OCT}
Do OCT là kỹ thuật tạo ảnh giao thoa sử dụng ánh sáng kết hợp thấp, ảnh thu được
thường chứa nhiễu speckle. Thành phần nhiễu này tạo ra các biến thiên cường độ có
dạng hạt, làm giảm độ tương phản giữa các lớp mô và có thể che khuất những cấu
trúc nhỏ. Vì tín hiệu nhiễu có thể ảnh hưởng đến đặc trưng mà mạng học sâu trích
xuất, luận văn khảo sát việc khử nhiễu trước khi huấn luyện thay vì mặc định rằng
mạng sẽ tự học cách bỏ qua speckle.

Non-Local Means (NLM) được Buades và cộng sự đề xuất dựa trên tính tự tương đồng
của ảnh: giá trị mới tại một điểm ảnh được tính bằng trung bình có trọng số của
các điểm ảnh có vùng lân cận tương tự, không chỉ của các điểm nằm gần về mặt
không gian \cite{buades2005}. Với ảnh đầu vào $I$, dạng rời rạc của phép lọc có
thể biểu diễn bởi
\begin{equation}
  \widehat{I}(p)=\frac{1}{Z(p)}
  \sum_{q\in\Omega_p}
  \exp\!\left(-\frac{\lVert P_p-P_q\rVert_2^2}{h^2}\right)I(q),
  \label{eq:nlm}
\end{equation}
trong đó $P_p$ và $P_q$ là hai patch lần lượt có tâm tại $p$ và $q$,
$\Omega_p$ là miền tìm kiếm, $h$ điều khiển mức suy giảm của trọng số và $Z(p)$
là hệ số chuẩn hóa.

\subsection{NLM trong khảo sát đối chứng}
NLM có khả năng làm trơn những vùng nhiễu trong khi tận dụng các patch lặp lại
để hạn chế làm mờ cấu trúc. Các biến thể NLM đã được nghiên cứu trực tiếp cho
khử speckle trong ảnh OCT \cite{aum2015}. Tuy nhiên, NLM nguyên bản được xây
dựng chủ yếu cho nhiễu cộng và việc lọc quá mạnh có thể loại bỏ đồng thời cả
tín hiệu bệnh lý. Vì vậy, NLM trong luận văn được xem là một giả thuyết tiền xử
lý cần kiểm chứng bằng đối chứng, không được mặc nhiên xem là cải thiện dữ liệu.
Bộ lọc được chọn cho cấu hình mới là Bilateral tại
Mục~\ref{sec:bilateral-selected}; NLM được giữ để đối chiếu kết quả trước đó.

\section{Vai trò của khảo sát khử nhiễu}
Khảo sát sáu bộ lọc trên 14 thể tích là cơ sở chính cho lựa chọn
Bilateral; khảo sát bảy bộ lọc trên một thể tích được giữ làm bằng
chứng bổ sung. Hai quy trình được phân biệt tại
Phụ lục~\ref{app:denoise-support}. Các lần chạy NLM trước đó
được giữ ở Mục~\ref{sec:sweep-3d-denoise-enc32} làm đối chứng lịch sử, không
được xem là kết quả của Bilateral. So sánh phân loại trước--sau còn lại
giữ cùng phân hoạch và kiến trúc nhưng không cùng ngân sách học;
giới hạn suy luận về hiệu quả bộ lọc được nêu tại Mục~\ref{sec:final-denoise-plan}.

\section{Thiết kế khảo sát các phương pháp khử nhiễu trên một khối ảnh OCT}
\label{sec:denoise-comparison-method}

\subsection{Mẫu khảo sát và quy trình xử lý thống nhất}
Khảo sát bổ sung sử dụng khối ảnh \texttt{data\_2404} từ Harvard-GF,
kích thước $200^3$, kiểu uint8, mang nhãn dương tính ($y=1$) theo nhật ký
của tác giả. Không có ảnh sạch chuẩn đi kèm. Số liệu được cung cấp từ lần chạy
\path{scripts/compare_denoise_methods.py}; số liệu này được ghi nhận
theo nguồn cung cấp, với giới hạn là thiếu đối chiếu độc lập với mã lệnh, CSV và siêu dữ liệu gốc.
Khảo sát chất lượng ảnh này tách biệt với kết quả phân loại của
\texttt{denoise-enc-32} và các mô hình Nhóm A.

Bảy bộ lọc được áp dụng hai chiều trên từng B-scan, lấy lát theo trục 0,
rồi ghép lại thành khối ảnh. Đầu ra được đưa về thang $[0,255]$ của ảnh gốc
để so sánh. Việc tái lập đòi hỏi xác định chính xác phép chuyển đổi, cắt ngưỡng và làm tròn về
uint8; co giãn min--max riêng từng đầu ra không phù hợp vì có thể thay đổi
các đại lượng đo. Với bộ lọc dùng đầu vào $[0,1]$, ảnh uint8 được chia cho
255 trước khi xử lý. Thời gian ghi trong nhật ký đo trên máy cá nhân
Windows, CPU 12 luồng, không phải thời gian GPU.

\subsection{Các bộ lọc và cấu hình khảo sát}
Bảng~\ref{tab:denoise-configs} ghi các tham số theo nguồn cung cấp. Khả năng
tái lập còn bị giới hạn bởi thông tin thiếu về phiên bản thư viện, chế độ xử lý biên
và số vòng lặp tối đa; các thông tin này cần được đối chiếu với cấu hình lần chạy,
không suy ra từ giá trị mặc định của phiên bản mới.

\begin{longtable}{@{}>{\raggedright\arraybackslash}p{0.22\textwidth}>{\raggedright\arraybackslash}p{0.39\textwidth}>{\raggedright\arraybackslash}p{0.31\textwidth}@{}}
  \caption{Các phương pháp và tham số trong khảo sát khối ảnh \texttt{data\_2404}}
  \label{tab:denoise-configs}\\
  \toprule
  \textbf{Phương pháp} & \textbf{Tham số theo nhật ký} & \textbf{Cách triển khai}\\
  \midrule\endfirsthead
  \multicolumn{3}{c}{\textit{Bảng~\thetable{} (tiếp theo)}}\\
  \toprule
  \textbf{Phương pháp} & \textbf{Tham số theo nhật ký} & \textbf{Cách triển khai}\\
  \midrule\endhead
  \bottomrule\endfoot
  Gaussian & $\sigma=1{,}5$ pixel & \path{skimage.filters.gaussian}\\
  Median & Cửa sổ $3\times3$ & \path{skimage.filters.median}\\
  Bilateral & \texttt{sigma\_color}=0,10; \texttt{sigma\_spatial}=4,0 & \path{denoise_bilateral}\\
  TV--Chambolle & \texttt{weight}=0,10 & \path{denoise_tv_chambolle}\\
  Wavelet & \texttt{db4}, ngưỡng mềm; \texttt{BayesShrink}; \texttt{rescale\_sigma=True} & \path{denoise_wavelet}\\
  NLM & $h=0{,}10$; patch 5; khoảng tìm kiếm 3; \texttt{fast\_mode=True} & \path{denoise_nl_means}\\
  BM3D & Profile \texttt{np}; \texttt{sigma\_psd}$\approx0{,}042$ & Gói \texttt{bm3d}\\
\end{longtable}

Gaussian và median là các đối chứng làm mịn đơn giản. Bilateral kết hợp
trọng số khoảng cách không gian và khác biệt cường độ
\cite{tomasi1998bilateral}. TV kết hợp điều chuẩn tổng biến phân với thành
phần khớp dữ liệu của mô hình ROF; khảo sát dùng thuật toán Chambolle
\cite{rudin1992tv,chambolle2004tv}. Tham số \texttt{weight}=0,10 được giữ
theo API, không đồng nhất với $\lambda$ trong mọi cách viết hàm mục tiêu;
trong tài liệu \texttt{denoise\_tv\_chambolle}, \texttt{weight} tương ứng
$1/\lambda$ với quy ước ROF được trình bày tại đó \cite{skimageRestoration}.

Wavelet dùng BayesShrink để chọn ngưỡng thích ứng cho các hệ số chi tiết
\cite{chang2000wavelet}. NLM dùng độ tương tự giữa các patch
\cite{buades2005}; với \texttt{patch\_distance}=3, tìm kiếm bị giới hạn
trong lân cận quy định, không phải so khớp trên toàn ảnh
\cite{skimageRestoration}. Việc coi cấu hình này trùng hoàn toàn với lần
tiền xử lý NLM trước cần đối chiếu nhật ký, vì cùng $h$ không đủ xác nhận cùng
thuật toán và tham số.

BM3D gom các khối 2D tương tự thành nhóm ba chiều để lọc cộng tác trong
miền biến đổi \cite{dabov2007bm3d}. Trong khảo sát này, thuật toán vẫn chạy
độc lập trên từng B-scan; chữ ``3D'' không có nghĩa là lọc trực tiếp toàn
khối ảnh OCT theo ba chiều không gian. Gói Python được dùng cung cấp giao diện gọi
các tệp thực thi nhị phân BM3D, không phải giao diện gọi OpenCV \cite{bm3dSoftware}.
Theo nhật ký, \texttt{sigma\_psd} lấy từ trung vị ước lượng trên 40 lát
bằng \texttt{estimate\_sigma}, ở thang $[0,1]$. Đây là tham số nhiễu
ước lượng để chạy bộ lọc, không phải phép đo trực tiếp nhiễu speckle thật.

\subsection{Vùng đo và các chỉ số đánh giá}
\label{sec:denoise-metric-definitions}
Nhật ký ghi dải tín hiệu theo chiều sâu $[29,57)$, nền phía trên
$[5,29)$, cắt 40 pixel ở biên ngang và dùng ngưỡng phân vị 60 của dải
ảnh gốc, tương ứng $84/255$. Theo bản rà soát quy trình do tác giả
cung cấp, chỉ ngưỡng cường độ được giữ chung; dải quanh đỉnh sáng
và các pixel thuộc vùng tín hiệu/nền được xác định lại trên ảnh sau
lọc. Vì vậy, các khoảng ghi trong nhật ký không xác nhận rằng mọi
phương pháp sử dụng hoàn toàn cùng vùng đo. Dải được gọi là RNFL
trong nhật ký không tương đương phân đoạn giải phẫu đã kiểm chứng. Trường
\texttt{dz=1}, thứ tự trục và cách dựng mặt nạ cần được đối chiếu
với mã đánh giá và siêu dữ liệu.

Gọi $\mu_s,\sigma_s$ là trung bình và độ lệch chuẩn trên vùng tín hiệu,
$\mu_b,\sigma_b$ là các đại lượng tương ứng trên nền. Với ảnh $I$ và vùng
$R$ gồm $N_R$ voxel, định nghĩa moment theo vùng là
\begin{equation}
 \mu_R=\frac1{N_R}\sum_{v\in R}I(v),\qquad
 \sigma_R=\sqrt{\frac1{N_R}\sum_{v\in R}[I(v)-\mu_R]^2}.
 \label{eq:denoise-roi-moments}
\end{equation}
Áp dụng lần lượt cho vùng tín hiệu và nền được xác định ở từng đầu ra
trong khảo sát ban đầu; không giả định các mặt nạ này trùng nhau.
Mẫu số $N_R$ mô tả độ phân tán của vùng đang đo; nếu mã nguồn dùng
$N_R-1$ để hiệu chỉnh phương sai mẫu thì phải ghi rõ và đối chiếu số liệu
đã lưu. Vùng rỗng hoặc mẫu số bằng 0 làm chỉ số tương ứng không xác định.
Theo định nghĩa thao tác của khảo sát,
\begin{align}
  \mathrm{SNR}_{\mathrm{ROI}}&=\frac{\mu_s}{\sigma_s},&
  \mathrm{ENL}&=\left(\frac{\mu_s}{\sigma_s}\right)^2,\\
  \mathrm{CNR}&=\frac{\mu_s-\mu_b}{\sqrt{\sigma_s^2+\sigma_b^2}}.
  \label{eq:denoise-roi-metrics}
\end{align}
SNR này là tỷ số moment trong vùng đo, không phải SNR công suất ở đơn vị
dB. ENL bằng bình phương SNR theo định nghĩa nên không phải một bằng chứng
độc lập thứ hai. Độ biến thiên trong vùng mô chứa cả cấu trúc lẫn nhiễu;
SNR hoặc ENL tăng có thể phản ánh làm phẳng cấu trúc thay vì chỉ giảm nhiễu.

Chỉ số biên $\beta$ là tương quan giữa độ lớn gradient của ảnh gốc và ảnh
sau lọc trên tập điểm biên mạnh được xác định từ ảnh gốc (nhóm 15\% cao
nhất theo mô tả). Gọi tập này là $\mathcal B$, đặt
$g_v=\|\nabla I_{\mathrm{goc}}(v)\|$ và
$h_v=\|\nabla I_{\mathrm{loc}}(v)\|$; khi đó
\begin{equation}
 \beta=\frac{\sum_{v\in\mathcal B}(g_v-\overline g)(h_v-\overline h)}
  {\sqrt{\sum_{v\in\mathcal B}(g_v-\overline g)^2}
  \sqrt{\sum_{v\in\mathcal B}(h_v-\overline h)^2}},\quad
 \overline g=\frac{\sum_{v\in\mathcal B}g_v}{|\mathcal B|},\quad
 \overline h=\frac{\sum_{v\in\mathcal B}h_v}{|\mathcal B|}.
 \label{eq:denoise-gradient-correlation}
\end{equation}
Khi xác định, $\beta\in[-1,1]$; gần 1 biểu thị tương quan thuận mạnh
giữa hai trường độ lớn gradient. SNR và ENL lớn biểu thị độ biến thiên
tương đối nhỏ trong vùng tín hiệu; CNR lớn biểu thị mức tách trung bình
tín hiệu--nền so với độ phân tán kết hợp. Công thức CNR ở đây giữ dấu,
không lấy trị tuyệt đối, phù hợp quy ước vùng tín hiệu sáng hơn nền.
Không có ngưỡng chung để một trong bốn chỉ số xác nhận ảnh tối ưu.
Việc tái lập phụ thuộc toán tử gradient, phạm vi chọn điểm biên, cách
gộp chỉ số theo lát/khối ảnh và quy ước xử lý phương sai bằng không. Đây là chỉ số
tham chiếu ảnh nhiễu gốc, không phải chỉ số hoàn toàn không tham chiếu.
Tương quan cao chỉ cho biết mức tương đồng của mẫu biến thiên gradient,
không bảo đảm bảo toàn độ lớn biên, tín hiệu bệnh lý hoặc RNFL.

\subsection{Tái lập và giới hạn lựa chọn tiền xử lý}
Kết quả được lưu trong thư mục \path{figures/denoise/}, gồm bảng
\path{denoise_metrics_2404.csv}, siêu dữ liệu
\path{denoise_compare_2404_all_meta.json} và cache đầu ra theo phương pháp.
Thời gian tính bộ lọc khi chưa có cache khác với thời gian đọc cache.
Việc đối chiếu chi phí đòi hỏi thông tin về CPU, RAM, số worker, số luồng mỗi worker,
phiên bản gói và phạm vi tính thời gian tải dữ liệu/I/O.
Các số đo này không so trực tiếp với khoảng 3,6 giờ tạo cache NLM đã nêu
cho một quy trình và môi trường khác.

Khảo sát một khối ảnh chỉ phục vụ thăm dò và không đủ để xác lập ưu
thế của bộ lọc. Lựa chọn Bilateral hiện tại dựa thêm trên khảo sát
14 thể tích ở Phụ lục~\ref{app:denoise-support}. Khối 2404 thuộc phân
hoạch kiểm tra; việc tham khảo mẫu này trong lựa chọn tiền xử lý
giới hạn tính độc lập của đánh giá trên chính phân hoạch đó. Đối
chứng phân loại ghép cặp là bằng chứng còn thiếu, thuộc nghiên cứu
ngoài phạm vi cuối cùng, không phải một lượt huấn luyện bổ sung
cần thực hiện. Theo Mục~\ref{sec:final-denoise-plan}, chỉ còn một
lượt tinh chỉnh 10 epoch trên ảnh khử nhiễu. SNR, ENL, CNR hay
$\beta$ riêng lẻ không xác nhận lợi ích phân loại của bộ lọc.
